% Source: https://artofproblemsolving.com/community/c3h2297889_proving_theorems
% Stable identity: aops-topic-2297889-q07
Let $B^*$ be a point of $m$ on the other side of $Q$ from $B$, so that $B^*$ and $A$ are on the same side of $t$.

\begin{asy}
usepackage("amsmath");
size(6.5cm);
defaultpen(fontsize(10pt));
pair Q=(0,0), P=(2/Tan(60),2);
pair dt=dir(60);
draw((-2.6,2)--(4.6,2),linewidth(0.9)); draw((-2.6,0)--(4.6,0),linewidth(0.9));
draw(Q-1.3*dt--P+1.3*dt,linewidth(0.9));
draw((2.4,2)--(2.6,2),Arrow(5)); draw((2.4,0)--(2.6,0),Arrow(5));
label("$\ell$",(4.6,2),E); label("$m$",(4.6,0),E); label("$t$",P+1.3*dt,N);
dot(P); dot(Q);
dot(P); label("$P$",P,NW);
dot(Q); label("$Q$",Q,SE);
dot(P+0.9*dt); label("$X$",P+0.9*dt,W);
dot(Q-0.9*dt); label("$Y$",Q-0.9*dt,E);
dot((4,2)); label("$A$",(4,2),N);
dot((-2,0)); label("$B$",(-2,0),S);
dot((3,0)); label("$B^{*}$",(3,0),S);
draw(arc(P,0.45,0,60),red); label("$\gamma$",P+0.77*dir(30.0),red+fontsize(8pt));
draw(arc(Q,0.45,180,240),blue); label("$\gamma$",Q+0.77*dir(210.0),blue+fontsize(8pt));
draw(arc(Q,0.45,0,60),deepgreen); label("$\gamma$",Q+0.77*dir(30.0),deepgreen+fontsize(8pt));
\end{asy}

\begin{itemize}
\item $\angle XPA$ and $\angle PQB^*$ are corresponding angles, so $\angle XPA\cong\angle PQB^*$.
\item The rays $QP$, $QY$ are opposite, and so are $QB^*$, $QB$. Hence $\angle PQB^*$ and $\angle YQB$ are vertical angles, and $\angle PQB^*\cong\angle YQB$.
\end{itemize}
Therefore $\angle XPA\cong\angle YQB$. $\blacksquare$
